\documentclass[11pt]{article}

\usepackage[a4paper,margin=1in]{geometry}
\usepackage{amsmath,amssymb,booktabs}
\usepackage[T1]{fontenc}
\usepackage{lmodern}
\usepackage{hyperref}

\newcommand{\dd}{\,\mathrm{d}}
\newcommand{\lap}{\nabla^2}
\newcommand{\Oh}{\mathcal{O}}
\newcommand{\Om}{\Omega}

\title{Quasi-Geostrophic Tutorial Test Cases}
\author{Pequod QG solver}
\date{\today}

\begin{document}
\maketitle

\section{Model and numerical conventions}

\paragraph{Unit conventions.}
The public \texttt{pequod.qg} and \texttt{pequod.qg2} interfaces now accept
SI inputs and expose SI field snapshots. Their numerical engines remain
dimensionless. The analytic benchmarks and legacy cases in
\texttt{tutorials/qg/main.py}, \texttt{convergence.py}, and
\texttt{tutorials/qg2/main.py} explicitly select
\texttt{units="dimensionless"}, preserving all benchmark values below.
The separate \texttt{wind\_experiment.py} runner uses SI inputs and a
final time in years, as specified in its subsection below.

The test cases implemented in \texttt{tutorials/qg/main.py} and
\texttt{tutorials/qg/convergence.py} use the one-layer
quasi-geostrophic (QG) potential-vorticity anomaly
\begin{equation}
  q = \lap\psi - \lambda^2\psi,
  \qquad
  \lambda^2 = L_D^{-2},
  \label{eq:pv}
\end{equation}
with velocity and Jacobian conventions
\begin{equation}
  (u,v)=(-\psi_y,\psi_x),
  \qquad
  J(\psi,q)=\psi_xq_y-\psi_yq_x.
  \label{eq:velocity-jacobian}
\end{equation}
The equation advanced by the solver can be written
\begin{equation}
  q_t + J(\psi,q) + \beta\psi_x
  = \nu\lap q-r\zeta+S,
  \qquad
  \zeta=\lap\psi=q+\lambda^2\psi,
  \label{eq:qg}
\end{equation}
where $S$ is the externally prescribed forcing, $\nu$ is the PV
diffusivity, and $r$ is the linear relative-vorticity drag coefficient.
The code represents the beta effect as the periodic term $\beta v=
\beta\psi_x$ rather than including the non-periodic function $\beta y$
inside $q$.

All solid boundaries are impermeable and free slip.  Thus
\begin{equation}
  \boldsymbol{u}\mathbin{\cdot}\boldsymbol{n}=0
  \quad\Longleftrightarrow\quad
  \psi=C_k
  \quad\text{on each connected solid boundary }\Gamma_k.
  \label{eq:solid-bc}
\end{equation}
The manufactured cases below choose every wall or island constant to be
zero.  The numerical PV flux normal to regular walls and obstacle faces is
also set to zero.  These cases therefore test a prescribed $C_k=0$
multiply-connected problem; they do not determine a general island
constant from an island rule or circulation constraint.

The exact manufactured cases initialize $q$ directly at conservative and
CABARET face points and supply the corresponding analytic velocity.  The
non-manufactured square-cylinder case instead initializes velocity and lets
the solver reconstruct a dynamically consistent PV field.

\section{Summary of cases in \texttt{main.py}}

The default resolution is $N_x=N_y=128$.  It can be changed with the
\texttt{--resolution} command-line option.

\begin{table}[ht]
  \centering
  \small
  \begin{tabular}{@{}llllll@{}}
    \toprule
    Case & Outer BC & Obstacle & $\beta$ & $\nu$ & $(\mathrm{CFL},T)$ \\
    \midrule
    Manufactured square hole & periodic $x,y$ & centered square
      & $0$ & $0$ & $(0.2,2)$ \\
    Periodic-beta square hole & periodic $x,y$ & centered square
      & $1$ & $0$ & $(0.2,2)$ \\
    Channel square island & periodic $x$, solid $y$ walls & centered square
      & $1$ & $0$ & $(0.2,2)$ \\
    Periodic square cylinder & periodic $x,y$ & centered square
      & $0$ & $10^{-3}$ & $(0.2,40)$ \\
    \bottomrule
  \end{tabular}
  \caption{Parameters selected in \texttt{tutorials/qg/main.py}.  Every
  case uses $L_D=1$.  The first three have $r=0$; the last has $r=0.01$.}
  \label{tab:cases}
\end{table}

\section{Doubly periodic manufactured square hole}
\label{sec:manufactured-hole}

This is the \texttt{manufactured\_square\_hole} case.  The fluid domain is
\begin{equation}
  \Om = [-\pi,\pi)\times[-\pi,\pi)\setminus\Oh,
  \qquad
  \Oh=[-a,a]\times[-a,a],
  \qquad a=\frac{\pi}{4}.
\end{equation}
The outer boundaries are periodic in both directions and
\begin{equation}
  \psi=0\qquad\text{on }\partial\Oh.
\end{equation}
Let
\begin{equation}
  A=1,
  \qquad L_D=1,
  \qquad c=\cos a,
  \qquad X=\cos x,
  \qquad Y=\cos y.
\end{equation}
The steady manufactured streamfunction is
\begin{equation}
  \psi_e=A(X-c)(Y-c).
  \label{eq:periodic-psi}
\end{equation}
It is exactly $2\pi$-periodic and vanishes on all four obstacle faces.  Its
PV and velocity are
\begin{align}
  q_e &=-A\left[
      X(Y-c)+Y(X-c)+\lambda^2(X-c)(Y-c)
    \right],
    \label{eq:periodic-q}\\
  u_e &= A(X-c)\sin y,
  &
  v_e &=-A\sin x\,(Y-c).
  \label{eq:periodic-velocity}
\end{align}
For $L_D=1$, equation~\eqref{eq:periodic-q} reduces to
\begin{equation}
  q_e=-A\left(3XY-2cX-2cY+c^2\right).
\end{equation}
The initial conditions are
\begin{equation}
  q(x,y,0)=q_e(x,y),
  \qquad
  \boldsymbol{u}(x,y,0)=\boldsymbol{u}_e(x,y).
\end{equation}
With $\beta=0$, $\nu=0$, and $r=0$, the forcing that makes the solution
steady is
\begin{equation}
  S_e=J(\psi_e,q_e)
  =A^2c\sin x\sin y\left(\cos x-\cos y\right).
  \label{eq:periodic-source}
\end{equation}
The order of the final cosine difference follows from the convention in
equation~\eqref{eq:velocity-jacobian}; reversing it gives $-J$.

The executable parameters are
\begin{equation}
  N_x=N_y=128,\quad
  \mathrm{CFL}=0.2,\quad
  T=2,\quad
  \nu=0,\quad
  \beta=0,\quad
  r=0.
\end{equation}

\section{Doubly periodic beta-plane square hole}
\label{sec:beta-hole}

The \texttt{periodic\_beta\_square\_hole} case has the same domain,
obstacle, boundary conditions, exact fields, and initial conditions as
Section~\ref{sec:manufactured-hole}.  Its parameters are
\begin{equation}
  A=1,\quad a=\frac{\pi}{4},\quad L_D=1,\quad
  \beta=1,\quad \nu=0,\quad r=0,\quad
  \mathrm{CFL}=0.2,\quad T=2.
\end{equation}
Because the beta effect is kept outside the definition of $q$, the
manufactured forcing is
\begin{align}
  S_{\beta,e}
    &=J(\psi_e,q_e)+\beta\psi_{e,x},\\
  \psi_{e,x}
    &=-A\sin x\,(\cos y-c),\\
  S_{\beta,e}
    &=A^2c\sin x\sin y(\cos x-\cos y)
      -\beta A\sin x(\cos y-c).
  \label{eq:beta-source}
\end{align}
Substitution into equation~\eqref{eq:qg} gives $q_t=0$, so this is again
an exact steady solution.

\section{Periodic channel with a square island}
\label{sec:channel-island}

This is the \texttt{channel\_square\_island} case.  Its outer rectangle is
\begin{equation}
  0\le x<L_x,
  \qquad
  -\frac{L_y}{2}\le y\le\frac{L_y}{2},
  \qquad
  L_x=2\pi,\quad L_y=\pi.
\end{equation}
The $x$ direction is periodic.  The north and south boundaries are
stationary impermeable free-slip walls.  The square island is centered at
$(x_0,0)$ and is
\begin{equation}
  \Oh=[x_0-a_x,x_0+a_x]\times[-a_y,a_y],
  \qquad
  x_0=\pi,\quad a_x=a_y=\frac{\pi}{4}.
\end{equation}
The boundary constants chosen by the manufactured solution are
\begin{equation}
  \psi=C_S=0\quad (y=-L_y/2),
  \qquad
  \psi=C_N=0\quad (y=L_y/2),
  \qquad
  \psi=C_I=0\quad (\partial\Oh).
  \label{eq:channel-bcs}
\end{equation}

Define
\begin{equation}
  k=\frac{2\pi}{L_x},
  \qquad
  \ell=\frac{2\pi}{L_y},
  \qquad
  \mu=\frac{\pi}{L_y},
\end{equation}
and the factors
\begin{align}
  f(x) &= \cos(k(x-x_0))-\cos(ka_x),\\
  w(y) &= \sin^2\!\left(\mu(y+L_y/2)\right),\\
  g(y) &= \cos(\ell y)-\cos(\ell a_y),\\
  h(y) &= w(y)g(y).
\end{align}
The exact steady streamfunction is
\begin{equation}
  \psi_e(x,y)=A f(x)h(y),
  \qquad A=1.
  \label{eq:channel-psi}
\end{equation}
The wall envelope $w$ vanishes at $y=\pm L_y/2$.  The factor $f$ vanishes
on the west and east island faces, while $g$ vanishes on the south and north
island faces.  Equation~\eqref{eq:channel-bcs} therefore holds exactly.

For compact expressions for the remaining fields, use
\begin{align}
  f'&=-k\sin(k(x-x_0)),
  &f''&=-k^2\cos(k(x-x_0)),
  &f'''&=k^3\sin(k(x-x_0)),\\
  w'&=\mu\sin(2\mu(y+L_y/2)),
  &w''&=2\mu^2\cos(2\mu(y+L_y/2)),
  &w'''&=-4\mu^3\sin(2\mu(y+L_y/2)),\\
  g'&=-\ell\sin(\ell y),
  &g''&=-\ell^2\cos(\ell y),
  &g'''&=\ell^3\sin(\ell y),
\end{align}
with
\begin{align}
  h' &= w'g+wg',\\
  h''&=w''g+2w'g'+wg'',\\
  h'''&=w'''g+3w''g'+3w'g''+wg'''.
\end{align}
The exact PV and velocity are
\begin{align}
  q_e &= A\left[f''h+f(h''-\lambda^2h)\right],
    \label{eq:channel-q}\\
  u_e &=-Af h',
  &v_e&=Af'h.
  \label{eq:channel-velocity}
\end{align}
For the exact forcing, define
\begin{align}
  q_{e,x}
    &=A\left[f'''h+f'(h''-\lambda^2h)\right],\\
  q_{e,y}
    &=A\left[f''h'+f(h'''-\lambda^2h')\right].
\end{align}
Then
\begin{equation}
  S_e
  =J(\psi_e,q_e)+\beta\psi_{e,x}
  =(Af'h)q_{e,y}-(Af h')q_{e,x}+\beta Af'h.
  \label{eq:channel-source}
\end{equation}
The initial conditions are $q(x,y,0)=q_e(x,y)$ and
$\boldsymbol{u}(x,y,0)=\boldsymbol{u}_e(x,y)$.  The remaining parameters are
\begin{equation}
  L_D=1,\quad \beta=1,\quad \nu=0,\quad r=0,\quad
  N_x=N_y=128,\quad \mathrm{CFL}=0.2,\quad T=2.
\end{equation}

This test exercises a multiply-connected Helmholtz inversion with a
prescribed island constant.  Determining a nonzero or time-dependent $C_I$
from circulation conservation is outside the scope of this manufactured
case.

\section{Doubly periodic square-cylinder vortex}
\label{sec:periodic-cylinder}

The \texttt{periodic\_square\_cylinder} case is a non-manufactured dynamic
test on
\begin{equation}
  [0,2\pi)\times[0,2\pi)
\end{equation}
with doubly periodic outer boundaries and a centered square obstacle of
width and height $\pi/2$.  Thus the obstacle is
\begin{equation}
  \Oh=[3\pi/4,5\pi/4]\times[3\pi/4,5\pi/4],
  \qquad \psi=0\quad\text{on }\partial\Oh.
\end{equation}
Let $A=0.35$, $\xi=x$, and $\eta=y$ for this $2\pi$-square.  The velocity
IVP is generated from
\begin{equation}
  \psi_{\mathrm{gen}}
  =A\left[\sin\xi+0.35\cos(2\xi+0.3)\right]\sin\eta,
\end{equation}
which gives
\begin{align}
  u_0
    &=-A\left[\sin\xi+0.35\cos(2\xi+0.3)\right]\cos\eta,\\
  v_0
    &=A\left[\cos\xi-0.70\sin(2\xi+0.3)\right]\sin\eta.
  \label{eq:cylinder-velocity}
\end{align}
The generating streamfunction does not itself satisfy the obstacle
condition.  During preprocessing, the solver takes the discrete curl of
equation~\eqref{eq:cylinder-velocity}, projects through the masked elliptic
solve so that $\psi=0$ in the square, and converts the resulting relative
vorticity to a consistent QG PV field.  Consequently this case has an
explicit velocity IVP but not a closed-form exact PV solution after the
obstacle projection.

There is no external forcing.  The parameters are
\begin{align}
  L_D&=1, & \beta&=0, & \nu&=10^{-3}, & r&=0.01,\\
  N_x&=N_y=128, & \mathrm{CFL}&=0.2, & T&=40.
\end{align}

\section{Convergence study in \texttt{convergence.py}}
\label{sec:convergence}

The convergence driver uses the f-plane manufactured square-hole problem
from Section~\ref{sec:manufactured-hole}, but its run parameters differ from
the interactive case:
\begin{equation}
  N\in\{64,128,256,512\},\qquad
  \mathrm{CFL}=0.15,\qquad
  T=0.5,\qquad
  \nu=\beta=r=0.
\end{equation}
The convergence controls are
\begin{itemize}
  \item \texttt{--resolutions} for the resolution list,
  \item \texttt{--final-time} for the final time, and
  \item \texttt{--cfl} for the CFL number.
\end{itemize}

Let $\Om_h$ be the active fluid-point mask.  On a uniform mesh the area
factor cancels from the relative discrete $L_2$ norm, so the reported error
is
\begin{equation}
  E_f^{(2)}
  =\left[
    \frac{\sum_{(i,j)\in\Om_h}(f_{h,ij}-f_{e,ij})^2}
         {\sum_{(i,j)\in\Om_h}f_{e,ij}^2}
  \right]^{1/2}.
  \label{eq:l2-error}
\end{equation}
The absolute maximum error is
\begin{equation}
  E_f^{(\infty)}
  =\max_{(i,j)\in\Om_h}|f_{h,ij}-f_{e,ij}|.
  \label{eq:linf-error}
\end{equation}
For successive resolutions $N_1$ and $N_2$, the observed order is
\begin{equation}
  p
  =\frac{\log(E_{N_1}/E_{N_2})}{\log(N_2/N_1)}.
  \label{eq:observed-order}
\end{equation}

The driver reports four relative-$L_2$ diagnostics:
\begin{enumerate}
  \item \textbf{PV inversion:} initialize with $q_e$, solve
    $(\lap-\lambda^2)\psi_h=q_e$, and report
    $E_{\psi,\mathrm{inv}}^{(2)}$ before time integration.
  \item \textbf{PV operator consistency:} set the analytic $\psi_e$ to zero
    on the solid mask, apply the periodic five-point discrete Helmholtz
    operator, and compare the resulting $q_h$ with $q_e$ to obtain
    $E_{q,\mathrm{op}}^{(2)}$.
  \item \textbf{Full streamfunction evolution:} integrate to $T=0.5$ with
    the exact source and report $E_{\psi,T}^{(2)}$.
  \item \textbf{Full PV evolution:} report $E_{q,T}^{(2)}$ at the same final
    time.
\end{enumerate}
It also records $E_{\psi,T}^{(\infty)}$, $E_{q,T}^{(\infty)}$, the maximum
absolute obstacle-normal numerical PV flux, and
\begin{equation}
  I_S=\sum_{(i,j)\in\Om_h}S_{e,ij}\,\Delta x\,\Delta y,
\end{equation}
which should vanish to roundoff by symmetry.  Results are printed as a
table and written to \texttt{qg\_manufactured\_convergence.csv} unless an
alternative path is supplied with \texttt{--output}.

\section{Commands}

From the repository root, the cases can be run with
\begin{verbatim}
PYTHONPATH=src .venv/bin/python tutorials/qg/main.py manufactured_square_hole
PYTHONPATH=src .venv/bin/python tutorials/qg/main.py periodic_beta_square_hole
PYTHONPATH=src .venv/bin/python tutorials/qg/main.py channel_square_island
PYTHONPATH=src .venv/bin/python tutorials/qg/main.py periodic_square_cylinder
PYTHONPATH=src .venv/bin/python tutorials/qg/convergence.py
\end{verbatim}

The $512^2$ masked sparse factorization in the default convergence sequence
can require substantial memory.  A shorter smoke test is, for example,
\begin{verbatim}
PYTHONPATH=src .venv/bin/python tutorials/qg/convergence.py \
  --resolutions 32 64 --final-time 0.1
\end{verbatim}

\section{Closed basin and northwest-corner bump}

Three additional cases use the closed unit basin $[0,1]^2$ with
$\psi=0$ on all four regular walls.  In the bump cases the solid is
\[
  \Oh=[0,1/4]\times[3/4,1].
\]
It meets both the west and north boundaries; its exposed east and south
faces also have $\psi=0$.  There is no independent island boundary constant
because this solid is connected to the outer wall.  The omitted north
endpoint is included in the solid mask when plotting.

The cases \texttt{basin\_manufactured} and
\texttt{basin\_northwest\_manufactured} use
\[
 P(s)=s^3(1-s)^3,\qquad
 \psi_e=A f(x)g(y),\qquad A=0.01.
\]
For the empty basin, $f$ and $g$ are both $P$ normalized by
$\max_{0\le s\le1}|P(s)|$.  For the bump, normalize each of
$P(x)(x-1/4)^3$ and $P(y)(y-3/4)^3$ separately by its maximum
absolute value on $[0,1]$ to obtain $f$ and $g$.
The role of these factors is explained below.
With $L_D=1$, the exact fields are
\begin{align}
 q_e&=A(f''g+fg''-fg),&
 u_e&=-Afg',&v_e&=Af'g,\\
 S_e&=u_e q_{e,x}+v_e q_{e,y}+\beta v_e.
\end{align}
Initialize $q=q_e$ and $(u,v)=(u_e,v_e)$.  The validation cases use
$\beta=1$, $\nu=r=0$, $\mathrm{CFL}=0.15$, and $T=0.1$.
The grid is $128\times128$ by default.  Automated tests compare
$32^2$ and $64^2$ for inversion and full evolution.

\subsection{Why choose a cubic boundary factor?}

The polynomial $P$ is a design choice for the manufactured solution,
not a requirement of the QG equations.  Impermeability requires only
constant streamfunction along a solid boundary.  For the chosen constant
zero, this means $\psi_e=0$; it does not by itself require zero tangential
velocity or zero PV.

The cubic factors impose additional conditions deliberately:
\[
 P(0)=P(1)=0,\qquad
 P'(0)=P'(1)=0,\qquad
 P''(0)=P''(1)=0.
\]
To see why the first two derivatives matter, consider a vertical wall at
$x=x_b$.  The separable streamfunction gives
\begin{align}
 \psi_e(x_b,y)&=A f(x_b)g(y),\\
 u_e(x_b,y)&=-A f(x_b)g'(y),\\
 v_e(x_b,y)&=A f'(x_b)g(y),\\
 q_e(x_b,y)&=A\left[f''(x_b)g(y)+f(x_b)g''(y)
                         -L_D^{-2}f(x_b)g(y)\right].
\end{align}
Here $u_e$ is the normal velocity and $v_e$ is the tangential velocity.
Therefore the conditions have distinct roles:
\begin{enumerate}
 \item $f(x_b)=0$ makes $\psi_e=0$ and the normal velocity $u_e=0$.
 \item $f'(x_b)=0$ additionally makes the tangential velocity $v_e=0$.
 \item $f''(x_b)=0$, together with $f(x_b)=0$, additionally makes $q_e=0$.
\end{enumerate}
At a horizontal wall the same argument applies to $g$, with the roles of
the velocity components exchanged.  Cubic zeros are a simple way to make
all three conditions hold simultaneously.

For comparison, a simple zero of $f$ generally leaves $f'$ nonzero, so it
allows tangential flow.  A double zero makes $f=f'=0$ but generally leaves
$f''$ nonzero, so boundary PV can remain nonzero.  A triple zero makes
$f=f'=f''=0$.  This is why the construction uses exponent three.

\subsection{How the factors enforce the northwest bump conditions}

Before normalization, the bump factors are
\[
 \widetilde f(x)=x^3(1-x)^3(x-1/4)^3,\qquad
 \widetilde g(y)=y^3(1-y)^3(y-3/4)^3.
\]
The endpoint factors enforce the conditions on the regular basin walls.
The additional zero at $x=1/4$ enforces them on the exposed east face of
the bump; the zero at $y=3/4$ enforces them on its exposed south face.
The factors also vanish along the extensions of those lines into the
fluid.  These are features of the selected exact solution, not additional
solid boundaries in the numerical geometry.

Normalization sets
\[
 f=\frac{\widetilde f}{\max_{s\in[0,1]}|\widetilde f(s)|},
 \qquad
 g=\frac{\widetilde g}{\max_{s\in[0,1]}|\widetilde g(s)|}.
\]
It changes the magnitude but preserves every zero and its multiplicity.
For the empty basin the same definition applies with
$\widetilde f=\widetilde g=P$.

\subsection{Implementation and validation limits}

In \texttt{tutorials/qg/main.py}, the function
\texttt{\_basin\_polynomials(bump)} constructs $P$ with
\begin{verbatim}
p = Polynomial([0., 1.])**3 * Polynomial([1., -1.])**3
\end{verbatim}
NumPy's \texttt{Polynomial} stores coefficients in increasing powers:
\texttt{Polynomial([0., 1.])} represents $s$, and
\texttt{Polynomial([1., -1.])} represents $1-s$.
For the bump, the function multiplies by
\begin{verbatim}
p *= Polynomial([-root, 1.])**3
\end{verbatim}
using \texttt{root=0.25} for $f$ and \texttt{root=0.75} for $g$.
It then normalizes the factors.
The function \texttt{basin\_manufactured\_fields} evaluates their
derivatives to construct $\psi_e$, $q_e$, the velocity, and the analytic
source.  The factory \texttt{make\_basin\_case} selects these fields as
initial data.

This construction is compatible with the current solver's practice of
zeroing stored PV on the south and west regular boundary rows and on
the solid obstacle mask.  The manufactured fields also have zero PV on
the north and east boundaries.  Consequently, these tests check inversion,
forcing, time evolution, and the corner geometry for a particularly
compatible exact solution.  They do not establish accuracy for general
wall data with nonzero PV or nonzero tangential velocity.

Although the manufactured velocity happens to vanish on every solid face,
this does not turn the solver into a general no-slip QG formulation:
zero tangential velocity is a property of this selected solution.
It is not an extra condition imposed on arbitrary solutions by the QG
Helmholtz inversion.  Further benchmarks with nonzero tangential wall flow
and nonzero boundary PV are needed to test those more general situations.

The case \texttt{basin\_northwest\_bump} uses the same initial fields and
geometry as the manufactured bump case, but sets $S=0$, $r=0.01$,
and $T=2$.  It is an evolving, unforced example with no claimed steady
exact solution; $\beta=1$, $\nu=0$, and $\mathrm{CFL}=0.15$.

Run these cases from the repository root:
\begin{verbatim}
PYTHONPATH=src .venv/bin/python tutorials/qg/main.py basin_manufactured
PYTHONPATH=src .venv/bin/python tutorials/qg/main.py \
  basin_northwest_manufactured
PYTHONPATH=src .venv/bin/python tutorials/qg/main.py basin_northwest_bump
\end{verbatim}
The existing \texttt{convergence.py} driver still targets the periodic
f-plane square-hole case.

\section{Two-layer QG: \texttt{qg2.py} and \texttt{tutorials/qg2/main.py}}

The class \texttt{TwoLayerQG}, exported as \texttt{pequod.qg2}, implements
a rigid-lid two-layer model. Layer 1 is the upper layer and layer 2 is the
lower layer. Its PV relations replace the single-layer relation:
\begin{align}
 q_1&=\lap\psi_1+F_1(\psi_2-\psi_1),\\
 q_2&=\lap\psi_2+F_2(\psi_1-\psi_2),\\
 \partial_t q_i+J(\psi_i,q_i)+\beta\partial_x\psi_i
   &=\nu_i\lap q_i-r_i\lap\psi_i+S_i,\qquad i=1,2.
\end{align}
The velocity convention is $(u_i,v_i)=(-\psi_{i,y},\psi_{i,x})$.
The positive stretching coefficients are conventionally
$F_i=f_0^2/(g'H_i)$, where $H_i$ is the reference layer depth and $g'$
is reduced gravity.  There is no additional single-layer
$-L_D^{-2}\psi_i$ term. The baroclinic deformation radius is
$(F_1+F_2)^{-1/2}$.

\subsection{Coupled inversion, boundaries and time stepping}

The inversion uses the weighted barotropic mode and the layer difference:
\begin{align}
 \psi_B&=\frac{F_2\psi_1+F_1\psi_2}{F_1+F_2},&
 q_B&=\frac{F_2q_1+F_1q_2}{F_1+F_2},&
 \lap\psi_B&=q_B,\\
 \psi_C&=\psi_1-\psi_2,&q_C&=q_1-q_2,&
 [\lap-(F_1+F_2)]\psi_C&=q_C.
\end{align}
Reconstruction gives
\[
 \psi_1=\psi_B+\frac{F_1}{F_1+F_2}\psi_C,\qquad
 \psi_2=\psi_B-\frac{F_2}{F_1+F_2}\psi_C.
\]
Both inverses use the five-point discrete Laplacian.  A completely
periodic domain without an obstacle requires zero mean $q_B$; the code
checks this compatibility condition and chooses zero mean $\psi_B$.
The constant baroclinic mode is retained.

Both layers share the same boundary geometry.  Periodic, channel and basin
boundaries, centered obstacles, and the existing basin corner bumps are
supported. By default each wall/obstacle has prescribed $\psi_i=0$, and the inherited
QG masks close the normal PV flux. Island constants are not inferred from
circulation. The limitations of the single-layer wall-PV masking described
above also apply to each layer.

The solver advances both CABARET transport engines with one common
timestep, restricted by both layers' advection and diffusion limits and
an explicit beta/drag cap. Both PV predictors are completed before a joint
inversion, and both correctors are completed before the next joint
inversion. Thus the velocities use a common temporal stage.
Drag is applied to relative vorticity $\lap\psi_i$, not directly to $q_i$.

\subsection{Manufactured two-layer cases}

Four steady benchmarks are selectable:
\begin{enumerate}
 \item \texttt{periodic\_square\_hole}: the periodic square-hole geometry and
   base streamfunction of Section~\ref{sec:manufactured-hole}.
 \item \texttt{channel\_square\_island\_manufactured}: the channel/island geometry and
   base streamfunction of Section~\ref{sec:channel-island}.
 \item \texttt{basin\_manufactured}: the empty unit basin with the
   normalized cubic factors described above.
 \item \texttt{basin\_northwest\_manufactured}: the unit basin with
   $\Oh=[0,1/4]\times[3/4,1]$ and the corresponding cubic bump factors.
\end{enumerate}
Let $\phi$ denote the selected base streamfunction and $z=\lap\phi$.
Its amplitude is 1 for the periodic and channel cases and 0.01 for the
basin cases.  Define $a_1=1$, $a_2=1/2$.  The exact steady layer fields are
\begin{align}
 \psi_{i,e}&=a_i\phi,&
 q_{i,e}&=a_i z+F_i(a_j-a_i)\phi,\qquad j=3-i,\\
 u_{i,e}&=-a_i\phi_y,&v_{i,e}&=a_i\phi_x,\\
 S_{i,e}&=a_i^2J(\phi,z)+\beta a_i\phi_x.
\end{align}
The stretching contribution cancels from the Jacobian because
$J(\phi,\phi)=0$. Both barotropic and baroclinic components are present,
although these proportional manufactured layer streamfunctions do not
test arbitrary nonlinear interactions between independent layer patterns.
Independent layer patterns and unequal coefficients are additionally
checked in the discrete inversion regression tests.

The initial PV and velocity are the exact fields above. All four cases use
\[
 F_1=1,\quad F_2=2,\quad \beta=1,\quad
 \nu_1=\nu_2=r_1=r_2=0,\quad
 \mathrm{CFL}=0.15,\quad T=0.1.
\]
The default grid is $64\times64$. Each case reports relative $L_2$ and
absolute $L_\infty$ errors for both $\psi_i$ and $q_i$ over active fluid
points. The option \texttt{--convergence} runs $N=32,64,128$ and reports
the measured orders for every error norm. These are full-evolution
errors; their observed rates need not equal the formal elliptic order.

\subsection{Two-layer square-island channel wake}

The name \texttt{channel\_square\_island} now selects a dynamic wake,
adapted from \texttt{square\_cylinder\_wake} in the Navier--Stokes tutorial.
The former exact channel benchmark has been renamed
\texttt{channel\_square\_island\_manufactured}.
The box is $[0,16)\times[0,8]$, periodic in $x$, with the solid
$[7.5,8.5]\times[3.5,4.5]$. The default CLI mesh is $256\times128$;
\texttt{--resolution N} means $N_y=N$, $N_x=2N$ for this case.
Both layers have stationary no-slip channel and island walls, enabled
explicitly with \texttt{no\_slip=True}; other benchmarks retain their
original impermeable QG boundary treatment.

The fixed bulk velocities are $(U_1,U_2)=(1,0.8)$, imposed by
\[
 \psi_{i,S}=0,\qquad \psi_{i,N}=-8U_i,\qquad \psi_{i,I}=-4U_i.
\]
Thus the cross-channel transport is $8U_i$ per layer even as the wake evolves.
These are prescribed island constants, not a circulation-constrained island rule.
Nonhomogeneous boundary lifts are included in both modal elliptic inversions;
adding a uniform velocity after inversion would violate island impermeability.
The no-slip wall PV is obtained from a one-sided Thom vorticity formula
plus the boundary stretching $F_i(\psi_j-\psi_i)$. Island-corner vorticity
uses the average of adjacent face closures. Dissipation remains $\nu\lap q_i$,
not an independently implemented momentum-viscosity operator.

Initial velocities are the same seeded Poiseuille profiles as the INCNS case:
with $\eta=y/8$, $k=2\pi/16$, $D=1$, and $\epsilon=10^{-3}$,
\[
 u_i=6U_i\eta(1-\eta)
       -\epsilon U_iD\,\sin(kx)\frac{\pi}{8}\sin(2\pi\eta),\qquad
 v_i=\epsilon U_iD\,k\cos(kx)\sin^2(\pi\eta).
\]
As in the INCNS startup these profiles do not initially satisfy the island
condition. Discrete velocity curls are inverted with the island/channel
Dirichlet data, then layer stretching is added to initialize PV consistently.
The evolved initial velocity is therefore the boundary-compatible projected
field, not the unmodified profile inside or across the obstacle.

Parameters are $F_1=0.2$, $F_2=0.05$ ($L_D=2$, depth ratio four),
$\beta=0$, $\nu_1=\nu_2=0.01$, $r_1=r_2=0$, and no external PV source.
The nominal $U_iD/\nu_i$ are 100 and 80. Constant transport represents
external driving; zero PV source does not mean an undriven decaying flow.
The final time is 120 and CFL is 0.15. The square has 16 cells per side
on the default mesh. This is a streamwise-periodic cylinder array, not
an isolated inlet/outlet wake, and shed vortices can recirculate through
the periodic seam. There is no analytic reference for this case; shedding
frequency and amplitude require time-series and grid-refinement checks.
\begin{verbatim}
PYTHONPATH=src .venv/bin/python tutorials/qg2/main.py \
  channel_square_island --resolution 128 --save-freq 150
\end{verbatim}

\subsection{Unforced two-layer northwest-bump example}

The unforced northwest-bump case uses the geometry and initial data of
the manufactured bump benchmark. It removes the
sustaining forcing ($S_1=S_2=0$), sets $(r_1,r_2)=(0,0.01)$,
and advances to $T=2$. Thus only the lower layer has direct linear drag;
the layers remain coupled through their streamfunctions. All other
parameters are as above. This evolving example has no claimed exact
reference solution.

\subsection{Wind-driven northwest-bump basin: idealized Gulf Stream}

The case \texttt{basin\_northwest\_bump\_wind} starts from rest,
$q_1=q_2=\psi_1=\psi_2=u_i=v_i=0$, on the unit basin with the
attached solid $[0,1/4]\times[1/2,1]$. Its exposed southeast corner
is at $(1/4,1/2)$, aligned with the wind-curl zero. Only this wind-driven
case uses the taller bump; the other northwest-bump cases are unchanged.
Both layers have $\psi_i=0$
on all regular walls and exposed bump faces (impermeability, not no-slip).
The time-independent wind acts only on the upper layer:
\[
 \frac{\tau_x}{\rho_0 H_1}=-\frac{W}{2\pi}\cos(2\pi y),\qquad
 \tau_y=0,\qquad S_1=-W\sin(2\pi y),\qquad S_2=0,\qquad W=0.1.
\]
These expressions are nondimensional. The implemented source is the
vertical wind-stress curl divided by upper-layer mass per unit area,
not the wind stress itself. The callback is
\texttt{northwest\_bump\_wind\_source} in \texttt{tutorials/qg2/main.py}.
The evolution equations and parameters are
\[
 (q_i)_t+J(\psi_i,q_i)+5(\psi_i)_x
 =5\times10^{-4}\nabla^2q_i-r_i\nabla^2\psi_i+S_i,
 \quad (r_1,r_2)=(0,0.05),\quad(F_1,F_2)=(40,10).
\]
The default is $N=64$, CFL $0.15$, and $T=20$. The coupling ratio
corresponds to $H_2/H_1=4$; the nondimensional baroclinic deformation
radius is $1/\sqrt{50}$. Dissipation is the solver's PV diffusion,
not a claim of a calibrated oceanic lateral-viscosity closure.

Negative southern wind curl drives an anticyclonic subtropical gyre;
positive northern curl drives a cyclonic subpolar gyre. Positive $\beta$
favors western boundary intensification, with a northward subtropical
western current and an eastward intergyre flow. The northwest bump
perturbs the northern circulation. This is an idealized Gulf Stream--like
spin-up experiment, not an Atlantic hindcast or a prescribed separation
at Cape Hatteras. Aligning the bump corner and wind-curl zero at $y=0.5$
provides a useful geometry for studying corner-associated separation,
but does not prescribe it: the current can turn along the solid face,
and its eventual detachment depends on the dynamics and resolution.
There is no exact reference or assertion of equilibrium at $T=20$.
Check longer spin-up and grid sensitivity before interpreting transport
or separation quantitatively; the nominal diffusion--beta scale
$(\nu/\beta)^{1/3}\simeq0.046$ spans only about three cells at $N=64$.
The removed northwest corner breaks the forcing's full-box zero-integral
symmetry; no artificial mean subtraction is applied to the wind curl.
For background on wind-driven western intensification, see
\url{https://ocw.mit.edu/courses/res-12-001-topics-in-fluid-dynamics-fall-2024/mitres_12_001_f24_essay3_pt4.pdf}.

\begin{verbatim}
PYTHONPATH=src .venv/bin/python tutorials/qg2/main.py \
  basin_northwest_bump_wind --resolution 64 --save-freq 50
\end{verbatim}
This writes the same four-panel, two-layer GIF as the other examples.
Use \texttt{--final-time} for longer spin-up and \texttt{--resolution 128}
for a better-resolved boundary current. Exact-error \texttt{--convergence}
is disabled for this dynamically forced case.

\paragraph{Spin-up and statistical equilibrium.}
Time is nondimensional, not days or years. $T=20$ is a short visualization
default, not an equilibrium integration. A long-wave baroclinic crossing
estimate is $T_R\sim L_x(F_1+F_2)/\beta=10$, while the direct lower-layer
drag time is $1/r_2=20$. Neither measures the slowest coupled adjustment.
For reference, unit-square PV diffusion has a lowest-mode decay time
$1/[\nu(\pi^2+\pi^2)]\simeq101$; the bump changes the spectrum.
As an initial experimental budget, use $T=500$ and extend to $1000$ or
beyond if statistics still drift. These values are timescale-based
guidance, not a demonstrated equilibration threshold.
Compare time-weighted block means and variability of each layer's kinetic
energy, total coupled energy, and gyre transport (streamfunction extrema)
over successive windows of length $50$--$100$. Discard initial spin-up;
require several successive windows without a systematic trend and with
mean differences small compared with sampling uncertainty (a few percent
is only a practical initial tolerance). Windows must exceed the measured
correlation time. Check stationarity of the mean jet path as well, then
repeat at finer resolution. A fluctuating statistically stationary state
need not approach a time-independent field.

\subsection{Wind experiment with SI inputs}

The reference length and velocity are fixed at
\[
 L_0=3\,840\,000\ \mathrm{m},\qquad U_0=0.01\ \mathrm{m\,s^{-1}},\qquad
 t_0=L_0/U_0=384\,000\,000\ \mathrm{s}\simeq12.17\ \mathrm{yr}.
\]
A year is defined as 365.25 days. Coordinates scale with $L_0$, time
with $t_0$, velocity with $U_0$, streamfunction with $U_0L_0$, PV and
relative vorticity with $1/t_0$, and PV sources with $1/t_0^2$.
The conversions for coefficients are
\[
 \beta^*=\beta L_0t_0,\qquad r^*=rt_0,\qquad
 \nu^*=\frac{\nu}{U_0L_0},\qquad F_i^*=F_iL_0^2,\qquad
 L_D^*=L_D/L_0.
\]
The physical defaults are
\[
 \beta=2\times10^{-11}\ \mathrm{m^{-1}s^{-1}},\quad
 r=4\times10^{-8}\ \mathrm{s^{-1}},\quad
 \nu=100\ \mathrm{m^2s^{-1}},
\]
giving $\beta^*=29491.2$, $r^*=15.36$, and $\nu^*=0.0026041667$.
Single-layer drag acts in that layer; two-layer drag defaults to $(0,r)$.
Public solver geometry/callback coordinates use metres, velocities m/s,
PV callbacks s$^{-1}$, and forcing callbacks s$^{-2}$. Public solver
\texttt{final\_time} is in seconds; only the wind CLI's
\texttt{--final-time} is in years. The numerical engine is available as
\texttt{solver.dimensionless}; modifying its fields uses model units.
Converted public field reads are snapshots, so IVP methods should be used
to prescribe initial state. Grid boundary tuples retain their dimensionless
selector; slip lengths are in metres and prescribed velocities in m/s.

The runner uses a physical $3840\times3840$ km basin, initially at rest.
Its default deformation radius is 38400 m on a $512^2$ grid: 5.12 spacings
per radius and a basin-to-radius ratio of 100. It sets
$F_1=0.8/L_D^2$, $F_2=0.2/L_D^2$, retaining the depth ratio four.
The default attached bump has width $0.25L_0$ and height $0.5L_0$.
With \texttt{--bump-cells K}, both dimensions are instead $K\Delta x$.
For independent dimensions, supply both \texttt{--bump-width-cells NX}
and \texttt{--bump-height-cells NY}; these set width $NX\Delta x$ and
height $NY\Delta y$, respectively. Do not combine this pair with
\texttt{--bump-cells}. For example, $NX=4$, $NY=8$ on the $128^2$ grid
gives a northwest bump 120 km wide and 240 km tall.

\paragraph{Wind stress and PV source.}
The generic solver adds a prescribed PV source $S_i$; it does not add
$\tau/(\rho_0H_1)$ directly to PV. The SI wind runner accepts a stationary
stress amplitude $W=\texttt{--wind-stress}$ and analytically computes
\[
 \tau_x=W\sin(\pi y/L),\quad\tau_y=0,\qquad L=3840\ \mathrm{km},
\]
\[
 S_1=\frac{\partial_x\tau_y-\partial_y\tau_x}{\rho_0H_1}
     =-\frac{\pi W}{\rho_0H_1L}\cos(\pi y/L),\qquad S_2=0.
\]
The default stress is 0.1 N/m$^2$, density 1025 kg/m$^3$, and upper-layer
depth 500 m, giving source amplitude $1.59633773\times10^{-13}$ s$^{-2}$.
Here $0\le y\le L$: the stress vanishes on both zonal walls and peaks
at midbasin. The curl is negative in the south and positive in the north.
This SI runner's stress profile differs from the older dimensionless
\texttt{basin\_northwest\_bump\_wind} case in \texttt{main.py}.
The density/depth normalize the wind forcing; coupling is specified by
$L_D$ and the fixed depth ratio. Alternatively \texttt{--wind W} directly
specifies the amplitude multiplying $-\cos(\pi y/L)$ in s$^{-2}$.
The options \texttt{--wind}
and \texttt{--wind-stress} are mutually exclusive.

\begin{verbatim}
PYTHONPATH=src .venv/bin/python -u tutorials/qg2/wind_experiment.py \
  --resolution 128 --radius 192000 --bump-cells 4 \
  --beta 2e-11 --drag 4e-8 --nu 100 \
  --wind-stress 0.1 --density 1025 --upper-depth 500 \
  --final-time 20 --interval 864000
\end{verbatim}
Here \texttt{--radius} is metres and \texttt{--interval} seconds (ten days).
The bump is $120\times120$ km at the northwest corner. The grid spacing
is 30 km and the radius has 6.4 spacings, but the nominal viscous--beta
width $(\nu/\beta)^{1/3}\simeq17.1$ km is less than one spacing.
Thus this coarse case is exploratory: it does not resolve that boundary
layer. The runner requires at least one spacing per deformation radius;
it no longer rejects refinement beyond ten spacings. This check alone
does not guarantee resolved boundary currents or meanders.
Old dimensionless CLI values must not be reused as SI inputs.

Each run creates a distinct \texttt{wind\_SI\_...} directory containing
\texttt{config.json}, compressed snapshots, \texttt{frames.json},
\texttt{diagnostics.jsonl}, and a four-panel GIF after completion.
The configuration records physical inputs, reference scales, and converted
dimensionless coefficients. For consistency with the numerical engine,
snapshot fields \texttt{q}, \texttt{psi}, and \texttt{t} remain dimensionless;
\texttt{t\_seconds}, \texttt{t\_years}, and \texttt{max\_velocity\_m\_s}
are also stored. Diagnostics label physical times and maximum speeds
explicitly, and the GIF displays metres, physical field units, and years.
Progress reports every 30 wall-clock seconds also show the maximum speed
in m/s. Snapshots are output data, not full restart checkpoints.

\paragraph{Obstacle-free rectangular basin: \texttt{wind\_experiment2.py}.}
This runner uses exactly the same SI parameters, initially resting two
layers, wind stress and upper-layer source as \texttt{wind\_experiment.py},
but removes the northwest bump entirely. The domain is the full
$[0,L]^2$ basin. All four regular walls retain $\psi_i=0$ and the same
Navier-slip closure; there is no irregular boundary. Both modal PV
inverses are plain sine transforms of the five-point operator, with no
capacitance factorization or obstacle Green-function correction.
The command-line options are unchanged except that all bump-size options
are absent. Outputs use the distinct prefix \texttt{wind2\_SI\_...}.

Both runners pass the physical source in s$^{-2}$ to the public solver.
With $q^*=T_0q$ and $t^*=t/T_0$, the tendency satisfies
$q_t=T_0^{-2}q^*_{t^*}$, so the numerical engine needs
\[
 S_i^*=T_0^2 S_i,\qquad T_0=L_0/U_0=384000000\ \mathrm{s}.
\]
\texttt{QGSIInterface.set\_forcing\_terms} applies this factor exactly
once, for both single- and two-layer solvers. The configuration's
\texttt{dimensionless.wind} is a record of the converted amplitude, not
an additional multiplication. For the default 0.1 N/m$^2$ stress,
the dimensionless source amplitude is approximately $23538.96$.

\subsection{Partial-slip lateral walls and runtime}

Both solvers accept \texttt{slip\_length} in metres (model length in
dimensionless mode), applied to every stationary regular wall and exposed
grid-aligned island/bump face. With normal $n$ directed into fluid,
\[
 u_n=0,\qquad u_t=\ell\,\partial_n u_t,\qquad\psi=C.
\]
Zero length gives no slip; infinite length gives zero tangential shear.
On a straight face, $\psi_n=\ell\zeta_w$ and $\psi_{nn}=\zeta_w$.
Taylor expansion from the wall to its fluid neighbour at distance $h$ gives
\[
 \psi(h)-C=\ell h\zeta_w+\tfrac12h^2\zeta_w+O(h^3),\qquad
 \zeta_w=\frac{2[\psi(h)-C]}{h(h+2\ell)}.
\]
Single-layer wall PV is $\zeta_w-L_D^{-2}C$; two-layer wall PV is
$\zeta_{i,w}+F_i(C_j-C_i)$. Obstacle corners average adjoining face
closures. Normal advective PV flux remains zero. Omitted north/east wall
nodes supply actual boundary PV in the diffusion stencil. Curved-wall
curvature terms are outside this grid-aligned rectangular geometry.
The explicit slip length overrides the old tangential wall-PV treatment;
regular wall tuples still specify stationary-wall data and impermeability.
The value \texttt{None} preserves the old manufactured-case mask, while
the old \texttt{no\_slip=True} channel option remains available.

The SI wind runner now defaults to \texttt{--slip-length 30000} on all
four outer walls and the exposed bump faces. Earlier wind runs used the
legacy free-slip-style mask, so their boundary physics differs.
For the completed $128^2$ run's geometry and physical coefficients:
\begin{verbatim}
PYTHONPATH=src MPLBACKEND=Agg .venv/bin/python -u \
  tutorials/qg2/wind_experiment.py \
  --resolution 128 --radius 180000 --bump-cells 4 \
  --wind-stress 0.1 --beta 2e-11 --drag 4e-8 --nu 100 \
  --slip-length 30000 --cfl 0.4 \
  --final-time 100 --interval 8640000
\end{verbatim}
This specifies 100 years and a 100-day snapshot interval; the bump is
$120\times120$ km. CFL 0.4 controls advection and explicit diffusion.
The explicit beta/drag source retains a separate restriction:
\[
 \Delta t\le\frac{\min(\mathrm{CFL},0.15)}
 {|\beta|\max(L_x,L_y)+\max_i r_i}.
\]
Increasing CFL therefore does not imply a proportional whole-run speedup.

A 1200-step $128^2$ profile found sparse triangular solves using 54\% of
stepping time, velocity reconstruction 8\%, and wall masking 7\%.
Symmetric minimum-degree sparse ordering, reusable modal/velocity arrays,
shared limiter increments, and cached wall indices/values reduce these
costs. A matched legacy-physics sample decreased from 6.65 s to 4.47 s;
the new partial-slip sample took 4.53 s. These are local profiled timings,
not a full-century timing claim. Details and a reproduction command are
in \texttt{tutorials/qg2/performance.md}.

\subsection{Two-layer API and execution}

PV callbacks return arrays with shape $(2,N_y,N_x)$. Velocity callbacks
return a pair $(u,v)$ with that leading layer axis on each component.
Use \texttt{set\_potential\_vorticity\_ivp(q,velocity)} to prescribe exact
PV, or \texttt{set\_ivp(velocity)} to derive PV from the sampled velocity
curls and coupled layer stretching. Time-independent forcing callbacks
also return arrays with a leading layer axis.

Examples from the repository root:
\begin{verbatim}
PYTHONPATH=src .venv/bin/python tutorials/qg2/main.py
PYTHONPATH=src .venv/bin/python tutorials/qg2/main.py \
  channel_square_island_manufactured --convergence
PYTHONPATH=src .venv/bin/python tutorials/qg2/main.py \
  basin_northwest_bump --resolution 128 --save-freq 50
\end{verbatim}
The CLI also accepts \texttt{--final-time}. By default, frames are saved
every five time steps; \texttt{--save-freq N} changes this interval.
The initial and final states are always included when saving is enabled.
NumPy \texttt{.npz} archives store
\texttt{t}, \texttt{q}, and \texttt{psi}; both field arrays include both
layers. Frames are written under
\begin{verbatim}
tutorials/qg2/data_<case>_<resolution>/
\end{verbatim}
Ordinary runs also generate
\texttt{tutorials/qg2/qg2\_<case>\_<resolution>.gif} with four panels:
rows correspond to layers 1 and 2, and columns show streamfunction and PV.
The time is displayed above the panels, and solid obstacles are shaded gray.
Color scales are fixed throughout the animation and shared between layers
for each variable. Only frames from the current run are used.
Use \texttt{--no-gif} to skip animation; adding \texttt{--save-freq 0}
also disables numerical frame output. Convergence runs skip GIF generation.
From Python, call \texttt{solver.solve(save\_freq=5)} followed by
\texttt{solver.save\_animation("qg2.gif", fps=12, dpi=100)}.
The single-layer \texttt{tutorials/qg/convergence.py} remains a separate
benchmark for the single-layer f-plane problem.

\end{document}
